\documentclass[11pt,a4paper]{article}

% ── Encodage & langue ────────────────────────────────────────────────────────
\usepackage[utf8]{inputenc}
\usepackage[T1]{fontenc}
% \usepackage[francais]{babel}  % desactive - non disponible dans cette installation

% ── Mise en page ─────────────────────────────────────────────────────────────
\usepackage[top=2.5cm, bottom=2.5cm, left=2.5cm, right=2.5cm]{geometry}
\usepackage{microtype}
\usepackage{parskip}

% ── Typographie ──────────────────────────────────────────────────────────────
\usepackage{palatino}
\usepackage{titlesec}
\usepackage{titletoc}

% ── Couleurs & boîtes ────────────────────────────────────────────────────────
\usepackage[dvipsnames,svgnames,x11names]{xcolor}
\usepackage{tcolorbox}
\tcbuselibrary{skins, breakable}

% ── Tableaux ─────────────────────────────────────────────────────────────────
\usepackage{booktabs}
\usepackage{tabularx}
\usepackage{longtable}
\usepackage{array}
\usepackage{multirow}

% ── Code & verbatim ──────────────────────────────────────────────────────────
\usepackage{listings}
\usepackage{fancyvrb}

% ── Liens ────────────────────────────────────────────────────────────────────
\usepackage{hyperref}
\hypersetup{colorlinks=true, linkcolor=AGTPrimary, urlcolor=AGTPrimary}

% ── En-têtes / pieds de page ─────────────────────────────────────────────────
\usepackage{fancyhdr}
\usepackage{lastpage}

% ── Icônes / symboles ────────────────────────────────────────────────────────
\usepackage{amssymb}
\usepackage{pifont}

% ── Palette de couleurs AGT ──────────────────────────────────────────────────
\definecolor{AGTPrimary}{HTML}{2563EB}
\definecolor{AGTDark}{HTML}{1E3A5F}
\definecolor{AGTLight}{HTML}{EFF6FF}
\definecolor{AGTGray}{HTML}{6B7280}
\definecolor{AGTBorder}{HTML}{E5E7EB}
\definecolor{AGTRed}{HTML}{DC2626}
\definecolor{AGTOrange}{HTML}{D97706}
\definecolor{AGTGreen}{HTML}{16A34A}
\definecolor{CodeBg}{HTML}{F8FAFC}

% ── Styles de sections ───────────────────────────────────────────────────────
\titleformat{\section}
  {\large\bfseries\color{AGTDark}}
  {\thesection}{1em}{}
  [\vspace{-0.3em}\color{AGTPrimary}\rule{\textwidth}{1.5pt}]

\titleformat{\subsection}
  {\normalsize\bfseries\color{AGTPrimary}}
  {\thesubsection}{1em}{}

\titleformat{\subsubsection}
  {\normalsize\bfseries\color{AGTDark}}
  {\thesubsubsection}{1em}{}

% ── En-têtes / pieds de page ─────────────────────────────────────────────────
\pagestyle{fancy}
\fancyhf{}
\fancyhead[L]{\small\color{AGTGray}\textbf{AGT Platform}}
\fancyhead[R]{\small\color{AGTGray}Chantier 1 — Conception B5 : Features 1--12}
\fancyfoot[C]{\small\color{AGTGray}page \thepage\ sur \pageref{LastPage}}
\renewcommand{\headrulewidth}{0.4pt}
\renewcommand{\footrulewidth}{0.4pt}

% ── Boîtes réutilisables ─────────────────────────────────────────────────────
\newtcolorbox{notebox}{
  breakable, enhanced,
  colback=AGTLight, colframe=AGTPrimary,
  boxrule=0.8pt, arc=4pt,
  left=8pt, right=8pt, top=6pt, bottom=6pt,
  fonttitle=\bfseries\small\color{AGTPrimary},
}

\newtcolorbox{warnbox}{
  breakable, enhanced,
  colback=Apricot!15, colframe=AGTOrange,
  boxrule=0.8pt, arc=4pt,
  left=8pt, right=8pt, top=6pt, bottom=6pt,
}

\newtcolorbox{featurebox}[1]{
  breakable, enhanced,
  colback=white, colframe=AGTBorder,
  boxrule=0.8pt, arc=6pt,
  left=10pt, right=10pt, top=8pt, bottom=8pt,
  title={\bfseries\color{AGTDark}#1},
  attach boxed title to top left={yshift=-2mm, xshift=5mm},
  boxed title style={colback=AGTPrimary, colframe=AGTPrimary, arc=4pt,
                     left=6pt, right=6pt, top=2pt, bottom=2pt},
  fonttitle=\bfseries\small\color{white},
}

\newtcolorbox{decisionbox}{
  breakable, enhanced,
  colback=gray!5, colframe=AGTGray,
  boxrule=0.5pt, arc=4pt,
  left=8pt, right=8pt, top=4pt, bottom=4pt,
}

% ── Listings code ────────────────────────────────────────────────────────────
\lstset{
  basicstyle=\ttfamily\small,
  backgroundcolor=\color{CodeBg},
  frame=single, framesep=4pt,
  rulecolor=\color{AGTBorder},
  breaklines=true,
  xleftmargin=10pt, xrightmargin=10pt,
  lineskip=1pt,
}

% ── Commandes utiles ─────────────────────────────────────────────────────────
\newcommand{\chantier}[1]{\textbf{\color{AGTPrimary}#1}}
\newcommand{\slug}[1]{\texttt{\color{AGTDark}#1}}
\newcommand{\todo}[1]{\textcolor{AGTRed}{\ding{46}~#1}}
\newcommand{\done}[1]{\textcolor{AGTGreen}{\ding{52}~#1}}
\newcommand{\warn}[1]{\textcolor{AGTOrange}{\textbf{⚠}~#1}}
\newcommand{\prio}[1]{\textcolor{AGTRed}{\textbf{#1}}}
\newcommand{\priom}[1]{\textcolor{AGTOrange}{\textbf{#1}}}
\newcommand{\priob}[1]{\textcolor{AGTGreen}{\textbf{#1}}}

% ─────────────────────────────────────────────────────────────────────────────
\begin{document}

% ── Page de garde ────────────────────────────────────────────────────────────
\begin{titlepage}
  \pagecolor{AGTDark}
  \color{white}
  \vspace*{3cm}
  \begin{center}
    {\Huge\bfseries AGT Platform}\\[0.6em]
    {\LARGE Chantier 1 — Finalisation \& Stabilisation}\\[2em]
    \color{AGTPrimary}
    \rule{0.6\textwidth}{2pt}\\[1.5em]
    \color{white}
    {\Large\bfseries Conception B5 : Base de connaissance \& Features}\\[0.5em]
    {\large Phase 0 — Audit et Conception des Modèles}\\[0.5em]
    {\large\color{AGTGray} Features 1 à 12 sur 24}\\[3em]
    \color{AGTGray}
    {\normalsize Document de référence — Session 32}\\[0.3em]
    {\normalsize Gabriel (Lead) — 18 mai 2026}\\[0.3em]
    {\normalsize AG Technologies — Confidentiel}
  \end{center}
  \vfill
  \begin{center}
    \color{AGTGray}{\small Suite : features 13 à 24 dans le prochain document}
  \end{center}
\end{titlepage}
\nopagecolor

% ── Table des matières ───────────────────────────────────────────────────────
\tableofcontents
\newpage

% ─────────────────────────────────────────────────────────────────────────────
\section{Préambule}
% ─────────────────────────────────────────────────────────────────────────────

Ce document est le résultat de l'audit approfondi des 12 premières features du bloc
\chantier{B5 (Chaîne métier complète)} du Chantier 1. Il sert de \textbf{boussole}
pour toutes les sessions qui travailleront sur B5.

\begin{notebox}
\textbf{Comment lire ce document.}
Chaque fiche feature décrit en français simple ce qui existe déjà, ce qui manque, et
exactement ce qu'il faut faire. Rien n'est ambigu. Aucune décision n'est laissée à
l'interprétation.
\end{notebox}

\begin{warnbox}
\textbf{Règle absolue :} aucune modification de modèle de données n'est faite sans
avoir relu la section correspondante dans ce document. Aucune initiative silencieuse.
Toute évolution non prévue ici est soumise à Gabriel avant action.
\end{warnbox}

% ─────────────────────────────────────────────────────────────────────────────
\section{Principes architecturaux validés}
% ─────────────────────────────────────────────────────────────────────────────

\subsection{Architecture en 3 couches}

Toute activité du bot laisse une trace dans l'une de ces 3 couches. Elles ne se
mélangent pas.

\begin{center}
\begin{tabularx}{\textwidth}{lllX}
\toprule
\textbf{Couche} & \textbf{Modèle} & \textbf{Qui écrit} & \textbf{Ce qu'elle capture} \\
\midrule
Communication & \slug{AIConversation / AIMessage} & Bot + Client & Le dialogue brut complet \\[4pt]
Exécution & \slug{AIActionLog} & Bot uniquement & Chaque action : type, payload, résultat, durée \\[4pt]
Résultat métier & \slug{Reservation, Commande...} & Bot + Entreprise & L'outcome business à valeur opérationnelle \\
\bottomrule
\end{tabularx}
\end{center}

\begin{notebox}
\textbf{Règle fondamentale :} toute feature active produit \textbf{au minimum une vue
stats} pour l'entreprise. Pour les features sans Résultat métier (FAQ, catalogues
d'orientation), la source des stats est \slug{AIActionLog}. Aucune feature n'est
laissée sans visibilité.
\end{notebox}

\subsection{Les 4 booléens sur le modèle Feature}

Ces 4 booléens pilotent le comportement du bot et de l'interface pour chaque feature.
Ils doivent être positionnés correctement dans le seeder.

\begin{center}
\begin{tabularx}{\textwidth}{lX}
\toprule
\textbf{Booléen} & \textbf{Signification} \\
\midrule
\slug{entreprise\_configure\_kb} & L'entreprise doit/peut configurer quelque chose dans \texttt{/knowledge} pour que cette feature fonctionne \\[4pt]
\slug{bot\_lit\_kb} & Le bot lit la KB de cette feature pendant la conversation avant d'agir \\[4pt]
\slug{bot\_ecrit\_result} & Le bot crée un enregistrement dans une table Résultat métier \\[4pt]
\slug{entreprise\_peut\_ecrire\_result} & L'entreprise peut intervenir sur les résultats (valider, annoter) avec mention de source \\
\bottomrule
\end{tabularx}
\end{center}

Sur le modèle \slug{ResultRecord} (base abstraite de tous les résultats) :
\begin{itemize}
  \item \slug{source} : \texttt{bot} / \texttt{entreprise} / \texttt{systeme}
  \item \slug{source\_user} : FK User nullable (si source = entreprise)
\end{itemize}

\subsection{Système de catalogues — Source canonique unique}

Trois systèmes coexistaient. \textbf{Décision définitive :} le Système 3 est la
source de vérité unique. Les Systèmes 1 et 2 sont dépréciés et leurs données migrées.

\begin{center}
\begin{tabularx}{\textwidth}{llX}
\toprule
\textbf{Système} & \textbf{Statut} & \textbf{Action} \\
\midrule
S1 : \slug{MenuCategorie/MenuPlat} & Déprécié & Migrer vers \slug{ItemCatalogue} + conserver le frontend \\[4pt]
S2 : \slug{CatalogueProduit}, \slug{CatalogueService}, etc. & Déprécié & Migrer + récupérer les champs spécialisés \\[4pt]
S3 : \slug{Catalogue/ItemCatalogue} & \textcolor{AGTGreen}{\textbf{Canonique}} & Enrichir avec les champs manquants \\
\bottomrule
\end{tabularx}
\end{center}

\begin{warnbox}
\textbf{Règle de migration :} on ne supprime rien tant que la migration n'est pas
validée par Gabriel sur un compte de test. Les anciens endpoints restent actifs
pendant la transition.
\end{warnbox}

\subsection{Règle frontend transversale}

\begin{notebox}
On n'expose sur le frontend \textbf{que ce qui fonctionne réellement côté backend}.
Aucun bouton qui mène nulle part. Les fonctionnalités non opérationnelles sont
masquées ou affichent \textit{"Bientôt disponible"}.
\end{notebox}

\subsection{Images des items — Plan V1 / V2}

\begin{center}
\begin{tabularx}{\textwidth}{lX}
\toprule
\textbf{Version} & \textbf{Approche} \\
\midrule
V1 — Chantier 1 & Bibliothèque d'images statiques préconfigurées par type d'item. L'entreprise choisit dans une liste. Stockées dans \texttt{public/}. \\[4pt]
V2 — Chantier 2 & Upload d'images réelles par l'entreprise — stockage cloud (S3/Cloudflare), miniatures automatiques. \\
\bottomrule
\end{tabularx}
\end{center}

% ─────────────────────────────────────────────────────────────────────────────
\section{Modifications du modèle de données}
% ─────────────────────────────────────────────────────────────────────────────

\subsection{Modèle \slug{Feature} — 4 nouveaux booléens}

\begin{lstlisting}
entreprise_configure_kb       : BooleanField  default=False
bot_lit_kb                    : BooleanField  default=False
bot_ecrit_result              : BooleanField  default=True
entreprise_peut_ecrire_result : BooleanField  default=False
\end{lstlisting}

Seeder à mettre à jour pour positionner ces booléens sur les 24 features.

\subsection{Modèle \slug{Ressource} — 11 nouveaux champs}

\begin{lstlisting}
duree_minutes            : PositiveIntegerField  default=30
delai_min_heures         : PositiveIntegerField  default=0
buffer_apres_minutes     : PositiveIntegerField  default=15
zone                     : CharField(50) choices=[terrasse, interieur, salon_vip, autre]
est_accessible_pmr       : BooleanField  default=False
tags                     : JSONField list  default=[]
chambre_type             : ForeignKey(ChambreType, null=True, blank=True)
trajet                   : ForeignKey(CatalogueTrajet, null=True, blank=True)
confirmation_automatique : BooleanField  default=False
politique_annulation     : CharField choices=[libre, avec_frais, non_remboursable]  default=libre
config_annulation        : JSONField  default={}
\end{lstlisting}

\begin{decisionbox}
\texttt{buffer\_slot\_min} sur \slug{ProfilEntreprise} est \textbf{déprécié}.
Conservé en base pour compatibilité descendante, non utilisé dans les nouveaux calculs.
\end{decisionbox}

\subsection{Modèle \slug{Reservation} — 6 nouveaux champs + statut}

\begin{lstlisting}
notif_client_envoyee     : BooleanField  default=False
notif_entreprise_envoyee : BooleanField  default=False
montant_rembourse        : DecimalField(10,2)  null=True, blank=True
date_remboursement       : DateTimeField  null=True, blank=True
motif_annulation         : TextField  blank=True
annule_par               : CharField choices=[client, entreprise, systeme]  blank=True
\end{lstlisting}

Ajouter le statut \texttt{en\_cours} (check-in effectué) pour les features
avec durée (chambre, billet).

\subsection{Modèle \slug{ChambreType} — 2 nouveaux champs}

\begin{lstlisting}
prix_nuit   : DecimalField(10,2)  null=True, blank=True
equipements : JSONField list  default=[]
\end{lstlisting}

\subsection{Modèle \slug{Commande} — 2 nouveaux champs}

\begin{lstlisting}
numero_commande : CharField(30)  unique=True  NOT NULL
                  Format : CMD-YYYYMMDD-XXXX  (genere a la creation)
notes_client    : TextField  blank=True
                  Visible par le bot au client. Distinct de notes (interne uniquement).
\end{lstlisting}

\subsection{Modèle \slug{ItemCatalogue} — 10 nouveaux champs}

\begin{lstlisting}
image_url                : URLField  blank=True
temps_preparation_min    : PositiveIntegerField  null=True, blank=True
allergenes               : JSONField list  default=[]
est_disponible_aujourd_hui : BooleanField  default=True
stock                    : PositiveIntegerField  null=True, blank=True
                           null = pas de gestion de stock (gere manuellement)
reference                : CharField(100)  blank=True
duree_min                : PositiveIntegerField  null=True, blank=True
type_service             : CharField(50)  blank=True
action_suivante          : CharField choices=[proposer_rdv, proposer_inscription,
                           proposer_devis, contact_direct, proposer_billet, aucune]
                           default=aucune
classes                  : JSONField  default=[]
                           Format : [{"nom": "Standard", "prix": 5000}, {"nom": "VIP", "prix": 8000}]
ville_depart             : CharField(100)  blank=True
ville_arrivee            : CharField(100)  blank=True
\end{lstlisting}

\subsection{Modèle \slug{QuestionFrequente} — 1 nouveau champ}

\begin{lstlisting}
ordre : PositiveIntegerField  default=0
\end{lstlisting}

\subsection{Nouveaux modèles à créer}

\subsubsection{\slug{ConsultationFAQ}}

\begin{lstlisting}
class ConsultationFAQ(models.Model):
    conversation      = FK(AIConversation)
    question_posee    = TextField()
    question_matchee  = FK(QuestionFrequente, null=True)  # null si aucune trouvee
    score_similarite  = FloatField(default=0)
    reponse_donnee    = TextField()
    source            = "bot"  # toujours bot
    created_at        = DateTimeField(auto_now_add=True)
\end{lstlisting}

\subsubsection{\slug{TransfertHumain}}

\begin{lstlisting}
class TransfertHumain(models.Model):
    STATUT = [en_attente, pris_en_charge, resolu]
    conversation      = FK(AIConversation)
    agence            = FK(Agence)
    motif             = TextField()
    statut            = CharField choices=STATUT  default=en_attente
    assigne_a         = FK(User, null=True)
    pris_en_charge_le = DateTimeField(null=True)
    resolu_le         = DateTimeField(null=True)
    source            = "bot"
    source_user       = FK(User, null=True)
    created_at        = DateTimeField(auto_now_add=True)
\end{lstlisting}

\subsection{Actions bot à créer}

\begin{center}
\begin{tabularx}{\textwidth}{llX}
\toprule
\textbf{Action} & \textbf{Paramètres} & \textbf{Description} \\
\midrule
\slug{search\_faq} & \texttt{query: str} & Recherche dans \slug{QuestionFrequente} par mots-clés. Remplace l'injection statique. Retourne top 3--5 résultats. \\[4pt]
\slug{get\_menu} & \texttt{categorie?} & Items actifs de \texttt{feature=menu\_digital}. \\[4pt]
\slug{get\_catalogue} & \texttt{query?, categorie?} & Items actifs filtrés. \\[4pt]
\slug{get\_services} & \texttt{query?} & Items services avec \texttt{type\_service} et \texttt{action\_suivante}. \\[4pt]
\slug{get\_trajets} & \texttt{depart?, arrivee?} & Filtre sur \texttt{ville\_depart}/\texttt{ville\_arrivee}. \\[4pt]
\slug{get\_room\_types} & aucun & \slug{ChambreType} actifs avec prix et équipements. \\[4pt]
\slug{get\_order\_status} & \texttt{numero\_commande} & Vérifie phone WhatsApp == contact.phone avant de retourner le statut. \\
\bottomrule
\end{tabularx}
\end{center}

% ─────────────────────────────────────────────────────────────────────────────
\section{Règles métier transversales}
% ─────────────────────────────────────────────────────────────────────────────

\subsection{Notifications après réservation ou commande}

S'applique à : \slug{prise\_rdv}, \slug{reservation\_table}, \slug{reservation\_chambre},
\slug{reservation\_billet}, \slug{orientation\_patient}, \slug{menu\_digital},
\slug{catalogue\_produits}.

\textbf{Déclencheur 1 — Création par le bot :} notifier l'entreprise + envoyer
un accusé de réception au client.

\textbf{Déclencheur 2 — Changement de statut par l'entreprise :} notifier le client
(confirmé / annulé).

Implémentation : signaux Django \texttt{post\_save}. Réutilise l'action
\slug{send\_email} existante. Champs \slug{notif\_client\_envoyee} et
\slug{notif\_entreprise\_envoyee} sur \slug{Reservation} pour éviter les doublons.

\subsection{Politique d'annulation par feature}

\begin{center}
\begin{tabularx}{\textwidth}{lll}
\toprule
\textbf{Feature} & \textbf{Politique par défaut} & \textbf{Configurable} \\
\midrule
\slug{prise\_rdv} & libre & Oui \\
\slug{reservation\_table} & libre & Oui \\
\slug{reservation\_chambre} & avec\_frais & Oui \\
\slug{reservation\_billet} & non\_remboursable & Oui \\
\slug{orientation\_patient} & libre & Oui \\
\bottomrule
\end{tabularx}
\end{center}

Remboursement = Transaction de type \texttt{credit} référençant la Reservation annulée.

\subsection{Livraison — Hors scope Chantier 1}

Statut \texttt{livree} existe en base mais masqué côté frontend.
Statuts visibles pour \slug{menu\_digital} et \slug{catalogue\_produits} :
\texttt{en\_attente / confirmee / en\_preparation / prete / annulee}.

\subsection{\slug{commande\_paiement} — Hors scope Chantier 1}

Requiert un provider externe opérationnel. Badge \textit{"Bientôt disponible"}
en Chantier 1. Implémentation en Chantier 2.

\subsection{FAQ — Injection statique à remplacer}

\begin{warnbox}
\textbf{Priorité haute.} Le \slug{ContextBuilder} injecte actuellement toute la FAQ
dans le system prompt (Bloc 2). Cela ruinera les tokens si une entreprise a 1 000
questions. \textbf{Action :} supprimer la FAQ du Bloc 2, créer l'action
\slug{search\_faq(query)} dynamique.
\end{warnbox}

\subsection{Sécurité suivi de commande}

Avant de partager tout statut de commande, le bot vérifie :
\texttt{Commande.contact.phone == conversation.contact.phone}.
En cas d'échec : TransfertHumain automatique.
\slug{numero\_commande} obligatoire, non nullable, généré à la création.

\subsection{Gestion du stock}

Le bot \textbf{lit} le stock pour vérifier la disponibilité.
Il ne le \textbf{modifie jamais}. Gestion exclusivement manuelle par l'entreprise.

\subsection{Confirmation automatique conditionnelle}

\slug{Ressource.confirmation\_automatique = True} ne confirme automatiquement
\textbf{que si le paiement est effectué}. Sans paiement : statut reste \texttt{en\_attente}.

% ─────────────────────────────────────────────────────────────────────────────
\section{Fiches features 1 à 12}
% ─────────────────────────────────────────────────────────────────────────────

% ─── Feature 1 ───────────────────────────────────────────────────────────────
\begin{featurebox}{Feature 1 — \slug{chatbot\_whatsapp}}

\textbf{Description métier :} Canal d'entrée WhatsApp. Le client envoie ses
messages, le bot répond. C'est la plomberie de tout le système.

\textbf{Ce qui existe :} \slug{WahaSession} (connexion QR, statut, numéro),
\slug{AIConversation}, \slug{AIMessage}, \slug{send\_text} via WAHA.

\medskip
\textbf{Booléens Feature :}
\begin{decisionbox}
\texttt{entreprise\_configure\_kb=False} | \texttt{bot\_lit\_kb=False} |
\texttt{bot\_ecrit\_result=True} (AIActionLog) | \texttt{entreprise\_peut\_ecrire\_result=False}
\end{decisionbox}

\medskip
\textbf{Actions à mener :}
\begin{enumerate}
  \item Vérifier que \slug{check\_quota("chatbot\_whatsapp")} est le \textbf{premier appel}
        du pipeline de traitement de tout message entrant.
  \item Créer le modèle \slug{TransfertHumain} (voir §2.8).
  \item Vérifier que chaque transfert humain crée un \slug{TransfertHumain} avec agence,
        motif, statut \texttt{en\_attente}.
\end{enumerate}

\end{featurebox}

\bigskip

% ─── Feature 2 ───────────────────────────────────────────────────────────────
\begin{featurebox}{Feature 2 — \slug{faq}}

\textbf{Description métier :} L'entreprise rédige ses questions/réponses
fréquentes. Le bot les consulte à la demande pour répondre aux clients.

\textbf{Ce qui existe :} \slug{FAQ} (conteneur), \slug{QuestionFrequente}
(questions avec catégorie), tab FAQ dans \texttt{/knowledge} — une seule FAQ
affichée (sous-exploité).

\medskip
\textbf{Booléens Feature :}
\begin{decisionbox}
\texttt{entreprise\_configure\_kb=True} | \texttt{bot\_lit\_kb=True} |
\texttt{bot\_ecrit\_result=True} (ConsultationFAQ) | \texttt{entreprise\_peut\_ecrire\_result=False}
\end{decisionbox}

\medskip
\textbf{Actions à mener :}
\begin{enumerate}
  \item Ajouter champ \slug{ordre} sur \slug{QuestionFrequente} (migration).
  \item Corriger l'interface \texttt{/knowledge} pour afficher plusieurs FAQ
        thématiques distinctes.
  \item \prio{Priorité haute :} supprimer l'injection statique de la FAQ dans le
        \slug{ContextBuilder} (Bloc 2). Créer action \slug{search\_faq(query)}.
  \item Créer modèle \slug{ConsultationFAQ} (voir §2.8). Logger chaque appel.
\end{enumerate}

\end{featurebox}

\bigskip

% ─── Feature 3 ───────────────────────────────────────────────────────────────
\begin{featurebox}{Feature 3 — \slug{prise\_rdv}}

\textbf{Description métier :} Un client prend rendez-vous via WhatsApp. Le bot
vérifie les disponibilités et crée le RDV automatiquement.

\textbf{Ce qui existe :} \slug{Ressource} (type=praticien/salle),
\slug{DisponibiliteRessource}, \slug{Reservation}, actions
\slug{check\_disponibilite} et \slug{create\_reservation}. Règle R14 active.

\medskip
\textbf{Booléens Feature :}
\begin{decisionbox}
\texttt{entreprise\_configure\_kb=True} | \texttt{bot\_lit\_kb=True} |
\texttt{bot\_ecrit\_result=True} | \texttt{entreprise\_peut\_ecrire\_result=True}
\end{decisionbox}

\medskip
\textbf{Actions à mener :}
\begin{enumerate}
  \item Ajouter sur \slug{Ressource} : \slug{duree\_minutes} (défaut 30),
        \slug{delai\_min\_heures} (défaut 0), \slug{buffer\_apres\_minutes} (défaut 15).
        Ces 3 champs profitent à toutes les features de réservation.
  \item Créer tab \textbf{"Disponibilités"} dans \texttt{/knowledge},
        visible si \slug{prise\_rdv} active.
  \item Mettre à jour \slug{check\_disponibilite} pour intégrer
        \slug{buffer\_apres\_minutes}.
  \item Page Résultats : filtrer par \texttt{feature.slug = "prise\_rdv"}.
  \item Boutons de changement de statut sur la page Résultats.
\end{enumerate}

\end{featurebox}

\bigskip

% ─── Feature 4 ───────────────────────────────────────────────────────────────
\begin{featurebox}{Feature 4 — \slug{reservation\_table}}

\textbf{Description métier :} Un client réserve une table au restaurant via
WhatsApp. Le bot vérifie et crée la réservation.

\textbf{Ce qui existe :} Même infrastructure que \slug{prise\_rdv}
(\slug{Ressource} type=table). Aucun tab frontend dédié.

\medskip
\textbf{Booléens Feature :}
\begin{decisionbox}
\texttt{entreprise\_configure\_kb=True} | \texttt{bot\_lit\_kb=True} |
\texttt{bot\_ecrit\_result=True} | \texttt{entreprise\_peut\_ecrire\_result=True}
\end{decisionbox}

\medskip
\textbf{Actions à mener :}
\begin{enumerate}
  \item Créer tab \textbf{"Tables"} dans \texttt{/knowledge},
        visible si \slug{reservation\_table} active.
  \item Ajouter sur \slug{Ressource} : \slug{zone}, \slug{est\_accessible\_pmr},
        \slug{tags}.
  \item Mêmes actions que \slug{prise\_rdv} pour buffer, Résultats, boutons.
\end{enumerate}

\end{featurebox}

\bigskip

% ─── Feature 5 ───────────────────────────────────────────────────────────────
\begin{featurebox}{Feature 5 — \slug{reservation\_chambre}}

\textbf{Description métier :} Un client réserve une chambre d'hôtel via WhatsApp.
Le bot présente les types disponibles et crée la réservation.

\textbf{Ce qui existe :} \slug{ChambreType} (catégories), \slug{Ressource}
(type=chambre), tab "Chambres" dans \texttt{/knowledge}. Pas d'interface pour
les chambres individuelles.

\medskip
\textbf{Actions à mener :}
\begin{enumerate}
  \item Ajouter sur \slug{ChambreType} : \slug{prix\_nuit}, \slug{equipements}.
  \item Ajouter sur \slug{Ressource} : \slug{chambre\_type FK}.
  \item Créer action bot \slug{get\_room\_types}.
  \item Créer interface pour les chambres individuelles dans \texttt{/knowledge}.
  \item Ajouter statut \texttt{en\_cours} sur \slug{Reservation}.
  \item Politique annulation : \texttt{avec\_frais} par défaut.
\end{enumerate}

\end{featurebox}

\bigskip

% ─── Feature 6 ───────────────────────────────────────────────────────────────
\begin{featurebox}{Feature 6 — \slug{reservation\_billet}}

\textbf{Description métier :} Un client réserve une place dans un transport via
WhatsApp. Le bot présente les trajets et émet le billet.

\textbf{Ce qui existe :} \slug{CatalogueTrajet} (lignes), \slug{Ressource}
(type=trajet), \slug{Reservation}.

\medskip
\textbf{Actions à mener :}
\begin{enumerate}
  \item Ajouter sur \slug{Ressource} : \slug{trajet FK}, \slug{confirmation\_automatique}
        (True pour transport — confirme si paiement reçu).
  \item Ajouter sur \slug{Reservation} : \slug{numero\_billet}, \slug{siege}.
  \item Politique annulation : \texttt{non\_remboursable} par défaut.
  \item Ajouter statut \texttt{en\_cours} (voyage en cours).
  \item Créer interface pour les véhicules dans \texttt{/knowledge}.
\end{enumerate}

\end{featurebox}

\bigskip

% ─── Feature 7 ───────────────────────────────────────────────────────────────
\begin{featurebox}{Feature 7 — \slug{menu\_digital}}

\textbf{Description métier :} Le restaurant publie son menu. Un client consulte
et commande directement via WhatsApp.

\textbf{Ce qui existe :} S1 (déprécié : \slug{MenuCategorie/MenuPlat}), S3
(canonique : \slug{ItemCatalogue}), \slug{Commande} + \slug{LigneCommande}.
Tab "Menu" branché sur S1.

\medskip
\textbf{Actions à mener :}
\begin{enumerate}
  \item \priom{Migration S1 → S3 :} rebrancher tab "Menu" sur \slug{ItemCatalogue}.
        Migrer les données \slug{MenuPlat} → \slug{ItemCatalogue}.
        Valider avec Gabriel sur un compte de test.
  \item Ajouter sur \slug{ItemCatalogue} : \slug{image\_url},
        \slug{temps\_preparation\_min}, \slug{allergenes},
        \slug{est\_disponible\_aujourd\_hui}.
  \item Images V1 : bibliothèque statique (§1.5).
  \item Créer action bot \slug{get\_menu(categorie?)}.
  \item Masquer statut \texttt{livree} côté frontend (§3.3).
\end{enumerate}

\end{featurebox}

\bigskip

% ─── Feature 8 ───────────────────────────────────────────────────────────────
\begin{featurebox}{Feature 8 — \slug{commande\_paiement} \textcolor{AGTRed}{— HORS SCOPE C1}}

\textbf{Description métier :} Passer et payer une commande directement via
WhatsApp avec un provider de paiement intégré.

\textbf{Décision :} Différée au Chantier 2. Requiert un provider externe
opérationnel (Orange Money, MTN MoMo).

\textbf{En Chantier 1 :} feature visible dans la marketplace, badge
\textit{"Bientôt disponible"}, non activable.

\end{featurebox}

\bigskip

% ─── Feature 9 ───────────────────────────────────────────────────────────────
\begin{featurebox}{Feature 9 — \slug{catalogue\_produits}}

\textbf{Description métier :} Une boutique publie son catalogue. Un client
consulte et commande via WhatsApp. Pas de paiement immédiat en C1.

\textbf{Ce qui existe :} S2 (déprécié : \slug{CatalogueProduit} avec \slug{stock},
\slug{reference}), S3 canonique sans ces champs.

\medskip
\textbf{Actions à mener :}
\begin{enumerate}
  \item \priom{Migration S2 → S3 :} récupérer \slug{stock} et \slug{reference}.
  \item Ajouter sur \slug{ItemCatalogue} : \slug{stock} (nullable), \slug{reference}.
  \item Le bot lit \slug{stock} uniquement. Ne le modifie \textbf{jamais} (§3.7).
  \item Créer action bot \slug{get\_catalogue(query?, categorie?)}.
\end{enumerate}

\end{featurebox}

\bigskip

% ─── Feature 10 ──────────────────────────────────────────────────────────────
\begin{featurebox}{Feature 10 — \slug{suivi\_commande}}

\textbf{Description métier :} Un client demande l'état de sa commande. Le bot
la retrouve et donne le statut en temps réel, sans intervention humaine.

\textbf{Ce qui existe :} \slug{Commande} avec statuts, \slug{LigneCommande},
\slug{Contact} (phone connu depuis WhatsApp).

\medskip
\textbf{Actions à mener :}
\begin{enumerate}
  \item Ajouter \slug{numero\_commande} : obligatoire, unique, format
        \texttt{CMD-YYYYMMDD-XXXX}, généré automatiquement.
  \item Ajouter \slug{notes\_client} : texte visible par le bot, distinct
        de \slug{notes} (interne).
  \item Créer action \slug{get\_order\_status(numero\_commande)} avec
        vérification : \texttt{Commande.contact.phone == conversation.contact.phone}.
        Échec → \slug{TransfertHumain} automatique.
\end{enumerate}

\end{featurebox}

\bigskip

% ─── Feature 11 ──────────────────────────────────────────────────────────────
\begin{featurebox}{Feature 11 — \slug{catalogue\_services}}

\textbf{Description métier :} Une entreprise de services liste ses prestations.
Le bot les présente et oriente vers la prochaine action selon la configuration.

\textbf{Ce qui existe :} S2 (déprécié : \slug{CatalogueService} avec
\slug{duree\_min}). Usage multi-secteurs : PME, Santé, Custom.

\medskip
\textbf{Actions à mener :}
\begin{enumerate}
  \item \priom{Migration S2 → S3 :} récupérer \slug{duree\_min}.
  \item Ajouter sur \slug{ItemCatalogue} : \slug{duree\_min},
        \slug{type\_service} (libre), \slug{action\_suivante} (choices).
  \item Créer action \slug{get\_services(query?)}. Le bot enchaîne selon
        \slug{action\_suivante} (ex : si \texttt{proposer\_rdv}, appelle
        \slug{check\_disponibilite}).
  \item Interface pour configurer \slug{type\_service} et \slug{action\_suivante}
        par prestation.
\end{enumerate}

\end{featurebox}

\bigskip

% ─── Feature 12 ──────────────────────────────────────────────────────────────
\begin{featurebox}{Feature 12 — \slug{catalogue\_trajets}}

\textbf{Description métier :} Une compagnie publie ses lignes, horaires et tarifs.
Le bot présente les trajets et oriente vers la réservation de billet.

\textbf{Ce qui existe :} S2 (déprécié : \slug{CatalogueTrajet} avec
\slug{horaires\_depart} JSONField). Feature complémentaire à
\slug{reservation\_billet}.

\textbf{Relation avec \slug{reservation\_billet} :}
\slug{CatalogueTrajet} = la ligne (horaires, tarifs).
\slug{Ressource} type=trajet = le véhicule physique.
Le FK \slug{Ressource.trajet → CatalogueTrajet} assure le lien.

\medskip
\textbf{Actions à mener :}
\begin{enumerate}
  \item \priom{Migration S2 → S3 :} récupérer les champs \slug{CatalogueTrajet}.
  \item Ajouter sur \slug{ItemCatalogue} : \slug{classes} (JSONField),
        \slug{ville\_depart}, \slug{ville\_arrivee}.
  \item \slug{action\_suivante = "proposer\_billet"} — fixe pour cette feature.
  \item Créer action \slug{get\_trajets(depart?, arrivee?)}.
\end{enumerate}

\end{featurebox}

% ─────────────────────────────────────────────────────────────────────────────
\section{Tableau récapitulatif des actions prioritaires}
% ─────────────────────────────────────────────────────────────────────────────

\begin{center}
\begin{longtable}{lp{7cm}p{3.5cm}l}
\toprule
\textbf{Priorité} & \textbf{Action} & \textbf{Feature(s)} & \textbf{Type} \\
\midrule
\endhead
\prio{Haute} & Corriger \slug{ContextBuilder} : supprimer FAQ Bloc 2, créer \slug{search\_faq} & FAQ & Backend \\[3pt]
\prio{Haute} & Ajouter \slug{numero\_commande} (non nullable, unique) & suivi\_commande & Migration \\[3pt]
\prio{Haute} & Ajouter les 4 booléens sur \slug{Feature} + seeder & Toutes & Migration \\[3pt]
\prio{Haute} & Ajouter les 11 champs sur \slug{Ressource} & rdv, table, chambre, billet & Migration \\[6pt]
\priom{Moyenne} & Migration S1/S2 → S3 (\slug{ItemCatalogue}) & menu, produits, services, trajets & Migration \\[3pt]
\priom{Moyenne} & Créer tab "Disponibilités" dans \texttt{/knowledge} & prise\_rdv & Frontend \\[3pt]
\priom{Moyenne} & Créer tab "Tables" dans \texttt{/knowledge} & reservation\_table & Frontend \\[3pt]
\priom{Moyenne} & Corriger tab FAQ multi-FAQ & faq & Frontend \\[3pt]
\priom{Moyenne} & Créer modèle \slug{TransfertHumain} & chatbot\_whatsapp & Backend \\[3pt]
\priom{Moyenne} & Créer modèle \slug{ConsultationFAQ} & faq & Backend \\[6pt]
\priob{Basse} & Ajouter 10 champs sur \slug{ItemCatalogue} & menu, produits, services, trajets & Migration \\[3pt]
\priob{Basse} & Créer les 7 actions bot & Multiples & Backend \\[3pt]
\priob{Basse} & Bibliothèque images statiques V1 & menu, produits & Frontend \\[3pt]
\priob{Basse} & Badge "Bientôt disponible" sur \slug{commande\_paiement} & commande\_paiement & Frontend \\
\bottomrule
\end{longtable}
\end{center}

% ─────────────────────────────────────────────────────────────────────────────
\section{Ordre d'exécution recommandé}
% ─────────────────────────────────────────────────────────────────────────────

\begin{notebox}
\textbf{Étape 1 — Migrations (faire en un seul lot).}
Feature (4 booléens + seeder) · Ressource (11 champs) · Reservation (6 champs +
statut en\_cours) · ChambreType (2 champs) · Commande (numero\_commande + notes\_client)
· ItemCatalogue (10 champs) · QuestionFrequente (ordre) · Créer TransfertHumain ·
Créer ConsultationFAQ.
\end{notebox}

\begin{notebox}
\textbf{Étape 2 — Backend actions.}
search\_faq (débloquer tests FAQ) · get\_menu · get\_catalogue · get\_services ·
get\_trajets · get\_room\_types · get\_order\_status (avec vérification sécurité).
\end{notebox}

\begin{notebox}
\textbf{Étape 3 — Migration données S1/S2 → S3.}
Valider sur compte de test avec Gabriel · Supprimer S1/S2 après validation.
\end{notebox}

\begin{notebox}
\textbf{Étape 4 — Frontend.}
Rebrancher tab Menu/Catalogue sur ItemCatalogue · Créer tabs Disponibilités + Tables ·
Corriger tab FAQ multi-FAQ · Badge Bientôt disponible sur commande\_paiement.
\end{notebox}

\vfill
\begin{center}
\color{AGTGray}
\rule{0.5\textwidth}{0.4pt}\\[0.5em]
{\small Document généré en session 32 — Gabriel (Lead) — 18 mai 2026}\\
{\small AG Technologies — Confidentiel}\\
{\small Suite : features 13 à 24 dans le prochain document de session}
\end{center}

\end{document}
